\chapter{Modèles discrets et suites affines}
\section*{Exercice 6.1 -- Déplacement d'un équilibre}
L'équilibre satisfait $\ell=0{,}75\ell+8$, soit $\ell=32$. La suite $v_n=u_n-32$ vérifie $v_{n+1}=0{,}75v_n$ et $v_0=-20$. Donc $u_n=32-20(0{,}75)^n$, qui tend vers $32$. L'apport de $8$ compense alors exactement la perte de $25\%$.

\section*{Exercice 6.2 -- Sans équilibre}
Ici $a=1$, donc $1-a=0$. La relation est arithmétique et donne $u_n=7+5n$. L'équation d'équilibre $\ell=\ell+5$ est impossible. Il faut traiter ce cas avant d'utiliser une formule contenant $1/(1-a)$.

\section*{Exercice 6.3 -- Alternance convergente}
L'équilibre est $\ell=2$. Ainsi $u_n=2-2(-1/2)^n$. Les termes d'indice pair sont inférieurs à $2$, les termes impairs supérieurs ; l'écart absolu vaut $2^{1-n}\to0$. La convergence n'exige pas la monotonie.

\section*{Exercice 6.4 -- Remboursement fictif}
En déroulant, $D_{12}=1000(1{,}01)^{12}-A[(1{,}01)^{12}-1]/0{,}01$. Imposer $D_{12}=0$ donne $A=10/[1-(1{,}01)^{-12}]$. Cette expression suppose le modèle exactement donné, sans frais ni arrondis mensuels ; elle n'est pas une description de tout contrat réel.

\section*{Exercice 6.5 -- Forçage variable}
Au rang $0$, $1-0-1=0$. Si $u_n=2^n-n-1$, alors $2u_n+n=2^{n+1}-n-2=2^{n+1}-(n+1)-1$. La récurrence prouve la formule. L'apport dépend ici du rang, donc le simple déplacement d'un équilibre constant n'est pas adapté.

\section*{Exercice 6.6 -- Tours de Hanoï}
Pour déplacer le plus grand disque, les $n$ petits doivent être sur l'autre tige, ce qui exige au moins $H_n$ coups. Il faut déplacer le grand puis reconstituer les petits, encore $H_n$ coups. Cette borne est atteinte par la procédure récursive. Ainsi $H_n+1$ est géométrique de raison $2$, donc $H_n=2^n-1$.
